\section*{Problem 2}

Set $f(x)=0$ to get the equilibria, then linearize at each one and check the eigenvalues of the Jacobian.

\textbf{System 1.} $\dot x_2=0$ gives $x_2=3x_1$, so $x_1(4-x_1^2)=0$ and $x_1\in\{0,\pm2\}$. With $J=\bmat{1-3x_1^2&1\\3&-1}$:
\begin{itemize}
\item $(0,0)$: $\lambda=\pm2$, saddle.
\item $(\pm2,\pm6)$: $\lambda^2+12\lambda+8=0$, $\lambda=-6\pm2\sqrt7$, stable node.
\end{itemize}

\textbf{System 2.} From $x_2(x_1-x_2)=0$: either $x_2=0$ (then $x_1=2$ or $-4$) or $x_2=x_1$ (then $x_1=2$ or $4$). $J=\bmat{2x_1-2x_2+2&-2x_1+4\\x_2&x_1-2x_2}$.
\begin{center}
\begin{tabular}{lll}
\toprule
equilibrium & eigenvalues & type\\
\midrule
$(2,0)$ & $6,\,2$ & unstable node\\
$(-4,0)$ & $-6,\,-4$ & stable node\\
$(2,2)$ & $2,\,-2$ & saddle\\
$(4,4)$ & $-1\pm j\sqrt7$ & stable focus\\
\bottomrule
\end{tabular}
\end{center}

\textbf{System 3.} $x_2(3-x_1-2x_2)=0$ and $x_1(2-x_1-x_2)=0$ give $(0,0)$, $(2,0)$, $(0,\tfrac32)$ and $(1,1)$. $J=\bmat{-x_2&3-x_1-4x_2\\2-2x_1-x_2&-x_1}$.
\begin{center}
\begin{tabular}{lll}
\toprule
equilibrium & eigenvalues & type\\
\midrule
$(0,0)$ & $\pm\sqrt6$ & saddle\\
$(2,0)$ & $-1\pm j$ & stable focus\\
$(0,\tfrac32)$ & $-\tfrac34\pm j\tfrac{\sqrt{15}}{4}$ & stable focus\\
$(1,1)$ & $-1\pm\sqrt2$ & saddle\\
\bottomrule
\end{tabular}
\end{center}

\textbf{System 4.} Let $a(x)=\tfrac12\tan(\pi x/2)$. The equilibrium conditions are $x_2=a(x_1)$ and $x_1=a(x_2)$. Subtracting them gives $x_2+a(x_2)=x_1+a(x_1)$, and since $x+a(x)$ is strictly increasing on $(-1,1)$ this forces $x_1=x_2=x$ with $\tan(\pi x/2)=2x$. The roots are $x=0,\pm\tfrac12$ (it is convex for $x>0$ with a negative slope at $0$, so there is only one positive root). The Jacobian is
\[
J=\bmat{-c&1\\1&-c},\qquad c=\tfrac{\pi}{4}\sec^2\!\big(\tfrac{\pi x}{2}\big),
\]
with eigenvalues $-c\pm1$.
\begin{itemize}
\item $(0,0)$: $c=\pi/4<1$, so $\lambda=-\tfrac\pi4\pm1=0.215,\,-1.785$, saddle.
\item $(\pm\tfrac12,\pm\tfrac12)$: $c=\pi/2$, so $\lambda=-0.571,\,-2.571$, stable node.
\end{itemize}
I only looked at $|x_i|<1$, the principal branch of $\tan$. The field has more equilibria on the other branches (for example $x_1=x_2\approx2.89$), and they are all stable nodes since $c>1$ there.

\begin{lstlisting}
sol = solve(f == 0, [x1 x2]);      % systems 1-3
xe  = double([sol.x1, sol.x2]);    % one equilibrium per row
A   = double(subs(J, x, xe(k,:).'));

g = @(t) t - 0.5*tan(pi*t/2);      % system 4, with x1 = x2 = t
r = fzero(g, [xs(k) xs(k+1)]);     % bracket from a sign change
\end{lstlisting}
